\documentclass[10pt,a4paper]{amsart}
\usepackage{amsmath,amssymb,amsthm,graphicx}
\usepackage{verbatim}
\newenvironment{reference}{\comment}{\endcomment}
%\newenvironment{reference}{}{}
\newenvironment{slogan}{\comment}{\endcomment}
\newenvironment{history}{\comment}{\endcomment}

% For commutative diagrams we use Xy-pic
\usepackage[all]{xy}

% We use 2cell for 2-commutative diagrams.
\xyoption{2cell}
% Theorem environments.
%
\theoremstyle{plain}
\newtheorem{theorem}[subsection]{Theorem}
\newtheorem{proposition}[subsection]{Proposition}
\newtheorem{lemma}[subsection]{Lemma}

\theoremstyle{definition}
\newtheorem{definition}[subsection]{Definition}
\newtheorem{example}[subsection]{Example}
\newtheorem{exercise}[subsection]{Exercise}
\newtheorem{situation}[subsection]{Situation}

\theoremstyle{remark}
\newtheorem{remark}[subsection]{Remark}
\newtheorem{remarks}[subsection]{Remarks}

\numberwithin{equation}{subsection}

% Macros
%
\def\lim{\mathop{\mathrm{lim}}\nolimits}
\def\colim{\mathop{\mathrm{colim}}\nolimits}
\def\Spec{\mathop{\mathrm{Spec}}}
\def\Hom{\mathop{\mathrm{Hom}}\nolimits}
\def\Ext{\mathop{\mathrm{Ext}}\nolimits}
\def\SheafHom{\mathop{\mathcal{H}\!\mathit{om}}\nolimits}
\def\SheafExt{\mathop{\mathcal{E}\!\mathit{xt}}\nolimits}
\def\Sch{\mathit{Sch}}
\def\Mor{\mathop{\mathrm{Mor}}\nolimits}
\def\Ob{\mathop{\mathrm{Ob}}\nolimits}
\def\Sh{\mathop{\mathit{Sh}}\nolimits}
\def\NL{\mathop{N\!L}\nolimits}
\def\CH{\mathop{\mathrm{CH}}\nolimits}
\def\proetale{{pro\text{-}\acute{e}tale}}
\def\etale{{\acute{e}tale}}
\def\QCoh{\mathit{QCoh}}
\def\Ker{\mathop{\mathrm{Ker}}}
\def\Im{\mathop{\mathrm{Im}}}
\def\Coker{\mathop{\mathrm{Coker}}}
\def\Coim{\mathop{\mathrm{Coim}}}

% Boxtimes
%
\DeclareMathSymbol{\boxtimes}{\mathbin}{AMSa}{"02}

%
% Macros for moduli stacks/spaces
%
\def\QCohstack{\mathcal{QC}\!\mathit{oh}}
\def\Cohstack{\mathcal{C}\!\mathit{oh}}
\def\Spacesstack{\mathcal{S}\!\mathit{paces}}
\def\Quotfunctor{\mathrm{Quot}}
\def\Hilbfunctor{\mathrm{Hilb}}
\def\Curvesstack{\mathcal{C}\!\mathit{urves}}
\def\Polarizedstack{\mathcal{P}\!\mathit{olarized}}
\def\Complexesstack{\mathcal{C}\!\mathit{omplexes}}
% \Pic is the operator that assigns to X its picard group, usage \Pic(X)
% \Picardstack_{X/B} denotes the Picard stack of X over B
% \Picardfunctor_{X/B} denotes the Picard functor of X over B
\def\Pic{\mathop{\mathrm{Pic}}\nolimits}
\def\Picardstack{\mathcal{P}\!\mathit{ic}}
\def\Picardfunctor{\mathrm{Pic}}
\def\Deformationcategory{\mathcal{D}\!\mathit{ef}}
\usepackage{cleveref}
\crefname{originallemma}{original lemma}{original lemmas}
\Crefname{originallemma}{Original lemma}{Original lemmas}
\makeatletter
\newlabel{lemma-duality-proper-over-field}{{0FVV}{source}{}{sourcepointer-1}{}}
\newlabel{lemma-duality-proper-over-field@cref}{{[originallemma][0][]0FVV}{[1][0][]source}{}{}{}}
\makeatother
\crefname{originallemma}{original lemma}{original lemmas}
\Crefname{originallemma}{Original lemma}{Original lemmas}
\makeatletter
\newlabel{coherent-lemma-coherent-filter}{{01YF}{source}{}{sourcepointer-2}{}}
\newlabel{coherent-lemma-coherent-filter@cref}{{[originallemma][0][]01YF}{[1][0][]source}{}{}{}}
\makeatother
\crefname{originallemma}{original lemma}{original lemmas}
\Crefname{originallemma}{Original lemma}{Original lemmas}
\makeatletter
\newlabel{lemma-lichtenbaum}{{0G5E}{source}{}{sourcepointer-3}{}}
\newlabel{lemma-lichtenbaum@cref}{{[originallemma][0][]0G5E}{[1][0][]source}{}{}{}}
\makeatother
\crefname{originallemma}{original lemma}{original lemmas}
\Crefname{originallemma}{Original lemma}{Original lemmas}
\makeatletter
\newlabel{properties-lemma-quasi-coherent-colimit-finite-type}{{01PG}{source}{}{sourcepointer-4}{}}
\newlabel{properties-lemma-quasi-coherent-colimit-finite-type@cref}{{[originallemma][0][]01PG}{[1][0][]source}{}{}{}}
\makeatother
\crefname{originallemma}{original lemma}{original lemmas}
\Crefname{originallemma}{Original lemma}{Original lemmas}
\makeatletter
\newlabel{cohomology-lemma-quasi-separated-cohomology-colimit}{{01FF}{source}{}{sourcepointer-5}{}}
\newlabel{cohomology-lemma-quasi-separated-cohomology-colimit@cref}{{[originallemma][0][]01FF}{[1][0][]source}{}{}{}}
\makeatother
\title{Lichtenbaum's theorem: vanishing of top cohomology is equivalent to no top dimensional component being proper}
\author{The Stacks Project Authors}
\date{stacks-project:duality.tex:local-snapshot-20261012; GNU FDL 1.2 or later}
\begin{document}
\maketitle
\noindent\textit{Editorial note.} All delivered mathematical text is the author's own, copied from the pinned Stacks Project chapter file duality.tex; no statement or proof step was rewritten, reordered or removed.\par\medskip
\noindent\textit{Editorial note.} Preamble provenance: the pinned chapter file opens with an input of the shared preamble and carries no preamble text of its own, so the project's single Stacks-wide preamble staged at sources/stacks-duality/preamble.tex is excerpted verbatim under the recipe fixed by LANE\_BRIEF 22.1, as three spans, lines 8-13 (verbatim package and the reference/slogan/history environments), lines 14-18 (Xy-pic setup) and lines 154-217 (the native theorem declarations and mathematical macros). Nothing was added, removed or redefined; the wrapper class and the amsmath/amssymb/amsthm/graphicx base packages are supplied by the assembler.\par\medskip
\noindent\textit{Editorial note.} The counted statement and proof contain no cross reference, no citation and no included file, so the delivered body is self contained.\par\medskip
\noindent\textit{Editorial note.} Canonical location of the counted proposition: Stacks tag 0G5F (sources/stacks-duality/pinned-original-source.tex line 9381).\par\medskip
\noindent\textit{Allowed standard prerequisites (editorial).}\begin{itemize}
\item The Lichtenbaum lemma bounding the dimension from nonvanishing top cohomology. Same chapter, Lemma labelled lemma-lichtenbaum, cited at proof lines 9424-9426. Turns nonvanishing of H\textasciicircum{}d on an integral subscheme into dimension d and properness.
\item Dualizing modules for a proper variety over a field. Same chapter, Lemma labelled lemma-duality-proper-over-field, cited at proof lines 9396-9400. Produces omega\_\{X'\} with H\textasciicircum{}d(X', omega\_\{X'\}) nonzero.
\item Filtrations of coherent modules with integral quotients. Cohomology of Schemes, Lemma labelled coherent-lemma-coherent-filter, cited at proof lines 9410-9412. Reduces the coherent case to a single integral closed subscheme with an ideal sheaf.
\item Quasi-coherent modules as filtered colimits of coherent submodules and vanishing of their cohomology. Properties, Lemma labelled properties-lemma-quasi-coherent-colimit-finite-type, cited at proof lines 9430-9434, and Cohomology, Lemma labelled cohomology-lemma-quasi-separated-cohomology-colimit, cited at proof lines 9434-9436. Proves the equivalence of (1) and (2).
\end{itemize}\medskip
\ifdefined\hypertarget\hypertarget{sourcepointer-1}{}\fi
\noindent\textit{Source reference.} Original lemma 0FVV: sources/stacks-duality/pinned-original-source.tex, source lines 6934--6934. Same-chapter pointer into the author environment lemma at source line 6934; invoked at proof line 9398 of the counted proof. It is retained only as a navigation pointer; every load-bearing step is inside the delivered span.\par\medskip
\ifdefined\hypertarget\hypertarget{sourcepointer-2}{}\fi
\noindent\textit{Source reference.} Original lemma 01YF: sources/stacks-chow/tags.txt, source lines 2554--2554. Cross-chapter standard prerequisite located in the project tag list staged in the local snapshot; invoked at proof line 9412 of the counted proof. It is retained only as a navigation pointer; every load-bearing step is inside the delivered span.\par\medskip
\ifdefined\hypertarget\hypertarget{sourcepointer-3}{}\fi
\noindent\textit{Source reference.} Original lemma 0G5E: sources/stacks-duality/pinned-original-source.tex, source lines 9329--9329. Same-chapter pointer into the author environment lemma at source line 9329; invoked at proof line 9424 of the counted proof. It is retained only as a navigation pointer; every load-bearing step is inside the delivered span.\par\medskip
\ifdefined\hypertarget\hypertarget{sourcepointer-4}{}\fi
\noindent\textit{Source reference.} Original lemma 01PG: sources/stacks-chow/tags.txt, source lines 2229--2229. Cross-chapter standard prerequisite located in the project tag list staged in the local snapshot; invoked at proof line 9434 of the counted proof. It is retained only as a navigation pointer; every load-bearing step is inside the delivered span.\par\medskip
\ifdefined\hypertarget\hypertarget{sourcepointer-5}{}\fi
\noindent\textit{Source reference.} Original lemma 01FF: sources/stacks-chow/tags.txt, source lines 1868--1868. Cross-chapter standard prerequisite located in the project tag list staged in the local snapshot; invoked at proof line 9436 of the counted proof. It is retained only as a navigation pointer; every load-bearing step is inside the delivered span.\par\medskip
\begin{theorem}
\label{theorem-lichtenbaum}
Let $X$ be a nonempty separated scheme of finite type over a field $k$.
Let $d = \dim(X)$. The following are equivalent
\begin{enumerate}
\item $H^d(X, \mathcal{F}) = 0$ for all coherent $\mathcal{O}_X$-modules
$\mathcal{F}$ on $X$,
\item $H^d(X, \mathcal{F}) = 0$ for all quasi-coherent $\mathcal{O}_X$-modules
$\mathcal{F}$ on $X$, and
\item no irreducible component $X' \subset X$ of dimension $d$
is proper over $k$.
\end{enumerate}
\end{theorem}

\begin{proof}
Assume there exists an irreducible component $X' \subset X$ (which we view
as an integral closed subscheme) which is proper and has dimension $d$.
Let $\omega_{X'}$ be a dualizing module of $X'$ over $k$, see
Lemma \ref{lemma-duality-proper-over-field}. Then
$H^d(X', \omega_{X'})$ is nonzero as it is dual to $H^0(X', \mathcal{O}_{X'})$
by the lemma. Hence we see that $H^d(X, \omega_{X'}) = H^d(X', \omega_{X'})$
is nonzero and we conclude that (1) does not hold.
In this way we see that (1) implies (3).

\medskip\noindent
Let us prove that (3) implies (1).
Let $\mathcal{F}$ be a coherent $\mathcal{O}_X$-module such that
$H^d(X, \mathcal{F})$ is nonzero. Choose a filtration
$$
0 = \mathcal{F}_0 \subset \mathcal{F}_1 \subset
\ldots \subset \mathcal{F}_m = \mathcal{F}
$$
as in Cohomology of Schemes, Lemma \ref{coherent-lemma-coherent-filter}.
We obtain exact sequences
$$
H^d(X, \mathcal{F}_i) \to H^d(X, \mathcal{F}_{i + 1}) \to
H^d(X, \mathcal{F}_{i + 1}/\mathcal{F}_i)
$$
Thus for some $i \in \{1, \ldots, m\}$ we find that
$H^d(X, \mathcal{F}_{i + 1}/\mathcal{F}_i)$ is nonzero.
By our choice of the filtration this means that there exists
an integral closed subscheme $Z \subset X$
and a nonzero coherent sheaf of ideals $\mathcal{I} \subset \mathcal{O}_Z$
such that $H^d(Z, \mathcal{I})$ is nonzero.
By Lemma \ref{lemma-lichtenbaum}
we conclude $\dim(Z) = d$ and $Z$ is proper over $k$
contradicting (3). Hence (3) implies (1).

\medskip\noindent
Finally, let us show that (1) and (2) are equivalent for any Noetherian scheme
$X$. Namely, (2) trivially implies (1). On the other hand, assume (1) and
let $\mathcal{F}$ be a quasi-coherent $\mathcal{O}_X$-module. Then we can write
$\mathcal{F} = \colim \mathcal{F}_i$ as the filtered colimit of its coherent
submodules, see
Properties, Lemma \ref{properties-lemma-quasi-coherent-colimit-finite-type}.
Then we have $H^d(X, \mathcal{F}) = \colim H^d(X, \mathcal{F}_i) = 0$
by Cohomology, Lemma \ref{cohomology-lemma-quasi-separated-cohomology-colimit}.
Thus (2) is true.
\end{proof}

\end{document}
